\documentclass[11pt]{beamer}

\usetheme{Boadilla}

%---------------------------------------------------------------------
% PACKAGES
%---------------------------------------------------------------------
\usepackage[english,greek]{babel}
\usepackage[utf8]{inputenc}
\usepackage[LGR,T1]{fontenc}
\usepackage{graphicx}
\usepackage{tikz}
\usepackage{xcolor}
\usepackage{booktabs}
\usepackage{listings}

% Shortcut για αγγλικό κείμενο μέσα σε ελληνικό
\newcommand{\en}[1]{\foreignlanguage{english}{#1}}

%---------------------------------------------------------------------
% FOOTLINE (όπως στο template)
%---------------------------------------------------------------------
\setbeamertemplate{footline}
{
  \leavevmode%
  \hbox{%
  \begin{beamercolorbox}[wd=.9\paperwidth,ht=2.25ex,dp=1ex,center]{author in head/foot}%
    \usebeamerfont{author in head/foot}%
    \insertsectionnavigationhorizontal{.7\paperwidth}{}{}%
  \end{beamercolorbox}%
  \begin{beamercolorbox}[wd=.1\paperwidth,ht=2.25ex,dp=1ex,right]{date in head/foot}%
    \usebeamerfont{date in head/foot}
    \insertframenumber{} / \inserttotalframenumber\hspace*{2ex}
  \end{beamercolorbox}}%
  \vskip0pt%
}
\beamertemplatenavigationsymbolsempty

%---------------------------------------------------------------------
% ΧΡΩΜΑΤΑ ΕΜΠ
%---------------------------------------------------------------------
\definecolor{NTUABlue}{RGB}{0,75,135}
\definecolor{Black}{RGB}{35,35,35}
\definecolor{White}{RGB}{255,255,255}
\definecolor{Darkgrey}{RGB}{98,100,104}
\definecolor{Middlegrey}{RGB}{192,193,195}
\definecolor{Lightgrey}{RGB}{230,230,231}

\setbeamercolor{item projected}{bg=NTUABlue,fg=White}
\setbeamercolor{section in toc}{fg=Black,bg=White}
\setbeamercolor{alerted text}{fg=NTUABlue}
\setbeamercolor*{palette primary}{fg=NTUABlue,bg=Lightgrey}
\setbeamercolor*{palette secondary}{fg=NTUABlue,bg=Lightgrey}
\setbeamercolor*{palette tertiary}{bg=NTUABlue,fg=White}
\setbeamercolor*{palette quaternary}{fg=NTUABlue,bg=Black}
\setbeamercolor{titlelike}{parent=palette primary,fg=NTUABlue}
\setbeamercolor{frametitle}{bg=Lightgrey, fg=NTUABlue}
\setbeamercolor{frametitle right}{bg=Middlegrey}

%---------------------------------------------------------------------
% LISTINGS STYLE
%---------------------------------------------------------------------
\lstset{
  basicstyle=\ttfamily\scriptsize,
  breaklines=true,
  frame=single,
  rulecolor=\color{Middlegrey},
  backgroundcolor=\color{Lightgrey!30},
  showstringspaces=false,
  language={},
}

%---------------------------------------------------------------------
% ΤΙΤΛΟΦΥΛΛΟ
%---------------------------------------------------------------------
\title[\en{Data Spaces} με \en{FIWARE/i4Trust}]{Περιβάλλον πειραματισμού με \en{Data Spaces} με χρήση \en{FIWARE} και \en{i4Trust}}

\author{Σωτήριος Κάκος}

\institute{Εθνικό Μετσόβιο Πολυτεχνείο \\ Σχολή Ηλεκτρολόγων Μηχανικών και Μηχανικών Υπολογιστών \\ Εργαστήριο Δικτύων Υπολογιστών}

\date{\today}

%=====================================================================
\begin{document}
%=====================================================================

% ─── Διαφάνεια 1: Τιτλοφύλλο ─────────────────────────────────────
\begin{frame}
\titlepage
\begin{center}
\small
\textbf{Επιβλέποντες:} Ε. Συκάς, Δ. Καλογεράς
\end{center}
\end{frame}

% ─── Διαφάνεια 2: Περιεχόμενα ────────────────────────────────────
\begin{frame}{Περιεχόμενα Παρουσίασης}
\tableofcontents
\end{frame}

%=====================================================================
\section{Εισαγωγή}
%=====================================================================

% ─── Διαφάνεια 3: Το πρόβλημα ────────────────────────────────────
\begin{frame}{Το Πρόβλημα του Διαμοιρασμού Δεδομένων}
Στη σύγχρονη ψηφιακή οικονομία, η ανταλλαγή δεδομένων μεταξύ οργανισμών
αποτελεί κρίσιμη παράμετρο για την καινοτομία και τη λήψη αποφάσεων.
Ωστόσο, ο διαμοιρασμός αυτός εγείρει σοβαρά ζητήματα:
\bigskip
\begin{itemize}
    \item Πώς διασφαλίζεται η εμπιστοσύνη μεταξύ οργανισμών που δεν
    γνωρίζονται μεταξύ τους;
    \item Πώς διατηρεί ο πάροχος τον έλεγχο επί των δεδομένων του
    μετά τη δημοσίευσή τους;
    \item Πώς επιτυγχάνεται η διαλειτουργικότητα χωρίς κεντρική
    αποθήκευση;
\end{itemize}
\bigskip
Οι παραδοσιακές προσεγγίσεις — κεντρικές πλατφόρμες ή διμερείς συμφωνίες —
δεν κλιμακώνονται και δημιουργούν εξάρτηση από ενδιάμεσους φορείς.
\end{frame}

% ─── Διαφάνεια 4: Data Spaces - η απάντηση ────────────────────────
\begin{frame}{Η Έννοια των \en{Data Spaces}}
Οι \en{Data Spaces} προτείνουν μια αποκεντρωμένη απάντηση στο πρόβλημα
του διαμοιρασμού δεδομένων.
\bigskip

\textbf{Βασικές αρχές:}
\begin{itemize}
    \item \textbf{Ψηφιακή Κυριαρχία:} Κάθε οργανισμός διατηρεί τον πλήρη
    έλεγχο επί των δεδομένων του.
    \item \textbf{Αποκεντρωμένη Αρχιτεκτονική:} Τα δεδομένα παραμένουν
    στις υποδομές του παρόχου.
    \item \textbf{Διαλειτουργικότητα:} Κοινά πρότυπα και πρωτόκολλα
    επιτρέπουν την επικοινωνία μεταξύ ετερογενών συστημάτων.
    \item \textbf{Πλαίσιο Εμπιστοσύνης:} Τεχνικοί μηχανισμοί
    (πιστοποιητικά, υπογραφές) αντικαθιστούν την εμπιστοσύνη σε
    μεσάζοντα.
\end{itemize}
\end{frame}

% ─── Διαφάνεια 5: Αντικείμενο και στόχοι ─────────────────────────
\begin{frame}{Αντικείμενο της Εργασίας}
Η παρούσα διπλωματική σχεδιάζει και υλοποιεί ένα λειτουργικό περιβάλλον
πειραματισμού — το \en{Plegma Data Space} — που αποδεικνύει στην πράξη
τον κύκλο εκχώρησης πρόσβασης σε δεδομένα.
\bigskip

\textbf{Συγκεκριμένα:}
\begin{itemize}
    \item Αξιοποίηση των ανοιχτών τεχνολογιών της \en{FIWARE Foundation}.
    \item Εφαρμογή του πλαισίου εμπιστοσύνης \en{i4Trust/iSHARE}.
    \item Μοντελοποίηση δεδομένων σε \en{NGSI-LD} με χρήση
    \en{Smart Data Models}.
    \item Επίδειξη ροών \en{Machine-to-Machine} τόσο για τον πάροχο όσο
    και για τον καταναλωτή.
    \item Τροφοδότηση του συστήματος με πραγματικά δεδομένα από το
    \en{Plegma Dataset}.
\end{itemize}
\end{frame}

%=====================================================================
\section{Θεωρητικό Υπόβαθρο}
%=====================================================================

% ─── Διαφάνεια 6: Δομικά στοιχεία Data Space ──────────────────────
\begin{frame}{Τα Δομικά Στοιχεία ενός \en{Data Space}}
Σύμφωνα με τη \en{Data Spaces Business Alliance (DSBA)}, ένας
\en{Data Space} αποτελείται από τέσσερις κατηγορίες δομικών στοιχείων:
\bigskip
\begin{itemize}
    \item \textbf{Διαχείριση Ταυτότητας \& Εμπιστοσύνης:} Επαλήθευση
    συμμετεχόντων και αξιολόγηση αξιοπιστίας.
    \item \textbf{Έλεγχος Πρόσβασης \& Χρήσης:} Πολιτικές που ορίζουν
    ποιος έχει πρόσβαση σε ποια δεδομένα.
    \item \textbf{Δημοσίευση \& Ανακάλυψη Δεδομένων:} Τυποποιημένη
    δημοσίευση πόρων και αναζήτηση από καταναλωτές.
    \item \textbf{Κυριαρχία Δεδομένων:} Διασφάλιση ότι κάθε
    συμμετέχων διατηρεί τον έλεγχο των δεδομένων του.
\end{itemize}
\bigskip
% TODO: εικόνα από το pdf της εκφώνησης (DSBA building blocks)
\begin{center}
\textit{[Εικόνα: \en{DSBA Building Blocks}]}
\end{center}
\end{frame}

% ─── Διαφάνεια 7: Πρωτοβουλίες ────────────────────────────────────
\begin{frame}{Διεθνείς Πρωτοβουλίες}
Η ανάπτυξη των \en{Data Spaces} υποστηρίζεται από διεθνείς οργανισμούς:
\bigskip
\begin{itemize}
    \item \textbf{\en{IDSA}}: \en{International Data Spaces Association} —
    το αναφορικό μοντέλο \en{IDS-RAM}.
    \item \textbf{\en{Gaia-X}}: Ευρωπαϊκή πρωτοβουλία για ομοσπονδιακή
    και ασφαλή ψηφιακή υποδομή.
    \item \textbf{\en{FIWARE Foundation}}: Ανοιχτά πρότυπα και λογισμικό,
    με ιδιαίτερη συνεισφορά στο πρότυπο \en{NGSI-LD}.
    \item \textbf{\en{DSBA}}: Κοινή πρωτοβουλία των παραπάνω για
    τεχνική σύγκλιση μεταξύ των οικοσυστημάτων.
\end{itemize}
\bigskip
Η παρούσα εργασία αξιοποιεί το πλαίσιο \en{i4Trust} — μια συνεργασία
\en{FIWARE} και \en{iSHARE} — που παρέχει συγκεκριμένη υλοποίηση των
δομικών στοιχείων με ανοιχτά εργαλεία.
\end{frame}

%=====================================================================
\section{Τεχνολογίες}
%=====================================================================

% ─── Διαφάνεια 8: PKI και X.509 ───────────────────────────────────
\begin{frame}{Κρυπτογραφία Δημόσιου Κλειδιού και Πιστοποιητικά \en{X.509}}
Η ταυτοποίηση των συμμετεχόντων βασίζεται στην κρυπτογραφία δημόσιου
κλειδιού. Κάθε οργανισμός διαθέτει ένα ζεύγος ιδιωτικού/δημόσιου
κλειδιού \en{RSA-2048}.
\bigskip

Τα ψηφιακά πιστοποιητικά \en{X.509} συνδέουν ένα δημόσιο κλειδί με μια
ταυτότητα, υπογεγραμμένα από μια Αρχή Πιστοποίησης (\en{CA}).
\bigskip

\textbf{Κρίσιμα στοιχεία της υλοποίησης:}
\begin{itemize}
    \item Το αναγνωριστικό \en{EORI} του οργανισμού καταχωρείται στο
    πεδίο \en{Serial Number} του πιστοποιητικού.
    \item Η επέκταση \texttt{keyUsage: digitalSignature} είναι
    υποχρεωτική για την επαλήθευση από τον \en{Keyrock}.
    \item Η αλυσίδα εμπιστοσύνης τερματίζεται σε μια τοπική \en{CA}
    (\en{PlegmaCA}).
\end{itemize}
\end{frame}

% ─── Διαφάνεια 9: JWT ─────────────────────────────────────────────
\begin{frame}[fragile]{\en{JSON Web Tokens} και η Επέκταση \en{iSHARE}}
Τα \en{JSON Web Tokens (JWT)} είναι αυτόνομα ψηφιακά \en{tokens} που
μεταφέρουν πληροφορίες υπογεγραμμένες ψηφιακά. Αποτελούνται από
τρία τμήματα: επικεφαλίδα, ωφέλιμο φορτίο, υπογραφή.
\bigskip

Το πλαίσιο \en{iSHARE} επεκτείνει τα \en{JWT} με δύο σημαντικές
προσθήκες:

\selectlanguage{english}
\begin{lstlisting}
Header:  { "alg": "RS256", "typ": "JWT",
           "x5c": ["<cert_leaf>", "<cert_CA>"] }

Payload: { "iss": "EU.EORI.NLPLEGMA",
           "sub": "EU.EORI.NLPLEGMA",
           "aud": "EU.EORI.NLPLEGMA",
           "iat": 1716000000,
           "exp": 1716000030 }
\end{lstlisting}
\selectlanguage{greek}

Η αλυσίδα πιστοποιητικών (\texttt{x5c}) επιτρέπει την αυτοτελή
επαλήθευση, ενώ η σύντομη διάρκεια ισχύος των 30 δευτερολέπτων
ελαχιστοποιεί τον κίνδυνο κατάχρησης.
\end{frame}

% ─── Διαφάνεια 10: iSHARE ─────────────────────────────────────────
\begin{frame}{Το Πλαίσιο Εμπιστοσύνης \en{iSHARE}}
Το \en{iSHARE} ορίζει τους κανόνες και τα πρωτόκολλα για ασφαλή
ανταλλαγή δεδομένων μεταξύ οργανισμών.
\bigskip

\textbf{Βασικοί ρόλοι:}
\begin{itemize}
    \item \textbf{\en{Service Provider}}: Διαχειρίζεται και
    δημοσιεύει τα δεδομένα.
    \item \textbf{\en{Service Consumer}}: Αιτείται πρόσβαση σε
    δεδομένα τρίτων.
    \item \textbf{\en{Satellite}}: Κεντρική αρχή εμπιστοσύνης —
    διατηρεί το μητρώο συμμετεχόντων.
    \item \textbf{\en{Authorization Registry (AR)}}: Αποθηκεύει
    και εκδίδει τις πολιτικές εκχώρησης.
\end{itemize}
\bigskip

Ο μηχανισμός \textbf{εκχώρησης εξουσιοδότησης (\en{delegation})}
επιτρέπει στον πάροχο να ορίζει με ακρίβεια τα δικαιώματα κάθε
καταναλωτή, εξασφαλίζοντας έτσι την ψηφιακή του κυριαρχία.
\end{frame}

% ─── Διαφάνεια 11: NGSI-LD ────────────────────────────────────────
\begin{frame}{Το Πρότυπο \en{NGSI-LD}}
Το \en{NGSI-LD} (\en{ETSI GS CIM 009}) είναι το πρότυπο \en{API} που
χρησιμοποιείται για την ανταλλαγή πλαισιακών πληροφοριών εντός του
\en{Data Space}.
\bigskip

\textbf{Μοντέλο πληροφορίας — τρεις θεμελιώδεις έννοιες:}
\begin{itemize}
    \item \textbf{Οντότητα (\en{Entity})}: Αντικείμενο με μοναδικό
    αναγνωριστικό \en{URI} και τύπο.
    \item \textbf{Ιδιότητα (\en{Property})}: Χαρακτηριστικό με τιμή,
    μονάδα μέτρησης και χρονοσφραγίδα.
    \item \textbf{Σχέση (\en{Relationship})}: Σύνδεση με άλλη οντότητα
    μέσω \en{URI}.
\end{itemize}
\bigskip

Το \texttt{@context} (από το \en{JSON-LD}) αντιστοιχίζει τα ονόματα
ιδιοτήτων σε παγκόσμια \en{URI}, δίνοντας στα δεδομένα σημασιολογική
βάση συμβατή με \en{Linked Data}.
\end{frame}

%=====================================================================
\section{Αρχιτεκτονική}
%=====================================================================

% ─── Διαφάνεια 12: Επισκόπηση αρχιτεκτονικής ──────────────────────
\begin{frame}{Αρχιτεκτονική του \en{Plegma Data Space}}
Η αρχιτεκτονική οργανώνεται σε \textbf{τρία λογικά επίπεδα}:
\bigskip
\begin{itemize}
    \item \textbf{Επίπεδο Επιβολής Πολιτικής}: \en{Kong API Gateway} +
    \en{DSBA-PDP}
    \item \textbf{Επίπεδο Εμπιστοσύνης \& Ταυτότητας}: \en{iSHARE
    Satellite} + \en{Keyrock}
    \item \textbf{Επίπεδο Δεδομένων}: \en{Orion-LD Context Broker} +
    \en{MongoDB}
\end{itemize}
\bigskip

% TODO: το component diagram εδώ
\begin{center}
\textit{[Εικόνα: \en{Component Diagram} του \en{Plegma Data Space}]}
\end{center}
\end{frame}

% ─── Διαφάνεια 13: Building blocks - Kong & PDP ───────────────────
\begin{frame}{Επίπεδο Επιβολής Πολιτικής}
Το επίπεδο αυτό αναχαιτίζει κάθε αίτημα πρόσβασης και αποφασίζει
αν θα προωθηθεί στα δεδομένα.
\bigskip

\textbf{\en{Kong API Gateway} — \en{Policy Enforcement Point (PEP)}:}
\begin{itemize}
    \item Μοναδικό σημείο εισόδου για όλες τις αιτήσεις (θύρα 8000).
    \item Αξιοποιεί το \en{plugin} \texttt{ngsi-ishare-policies}
    ειδικά σχεδιασμένο για \en{Data Spaces}.
    \item Λειτουργεί σε \en{DB-less} κατάσταση, με δηλωτική
    παραμετροποίηση μέσω \texttt{kong.yml}.
\end{itemize}
\bigskip

\textbf{\en{DSBA-PDP} — \en{Policy Decision Point (PDP)}:}
\begin{itemize}
    \item Αξιολογεί αιτήματα έναντι πολιτικών \en{iSHARE}.
    \item Επιστρέφει αποφάσεις \en{Permit} ή \en{Deny}.
    \item Υποστηρίζει σύνθετα σενάρια αξιολόγησης.
\end{itemize}
\end{frame}

% ─── Διαφάνεια 14: Trust & Identity layer ─────────────────────────
\begin{frame}{Επίπεδο Εμπιστοσύνης και Ταυτότητας}
\textbf{\en{iSHARE Satellite} — \en{Trust Anchor}:}
\begin{itemize}
    \item Διατηρεί το μητρώο εγκεκριμένων συμμετεχόντων.
    \item Παρέχει τη λίστα εμπιστευμένων \en{CA} (\texttt{trusted\_list}).
    \item Επαληθεύει την κατάσταση συμμετεχόντων μέσω του
    \texttt{/parties} \en{endpoint}.
\end{itemize}
\bigskip

\textbf{\en{Keyrock Identity Manager} — \en{Authorization Registry}:}
\begin{itemize}
    \item Διπλός ρόλος: πάροχος ταυτότητας \textit{και} μητρώο
    εξουσιοδότησης.
    \item Δέχεται \en{iSHARE JWT} και εκδίδει \en{access tokens}.
    \item Διαχειρίζεται τις πολιτικές εκχώρησης μέσω της βάσης
    \en{MySQL}.
    \item Μηχανισμός \en{extparticipant} για αυτόματη καταχώρηση
    εξωτερικών οργανισμών.
\end{itemize}
\end{frame}

% ─── Διαφάνεια 15: Data layer ─────────────────────────────────────
\begin{frame}{Επίπεδο Δεδομένων}
\textbf{\en{Orion-LD Context Broker}:}
\begin{itemize}
    \item Υλοποίηση \en{NGSI-LD} από τη \en{FIWARE Foundation}.
    \item Αποθηκεύει οντότητες σε μορφή \en{JSON-LD} στη
    \en{MongoDB}.
    \item Υποστηρίζει χρονικά ερωτήματα και ειδοποιήσεις αλλαγών.
\end{itemize}
\bigskip

\textbf{Σημαντική αρχιτεκτονική επιλογή:} Ο \en{Orion-LD}
εκτίθεται \textbf{αποκλειστικά} στο εσωτερικό δίκτυο
\en{Docker}. Δεν είναι προσβάσιμος εξωτερικά.
\bigskip

Έτσι εξασφαλίζεται ότι \textbf{καμία αίτηση δεν μπορεί να
παρακάμψει το \en{Kong} και τους μηχανισμούς εξουσιοδότησης}.
\end{frame}

%=====================================================================
\section{Μοντέλο Δεδομένων}
%=====================================================================

% ─── Διαφάνεια 16: Plegma Dataset ────────────────────────────────
\begin{frame}{Το \en{Plegma Dataset}}
Πραγματικά δεδομένα από 13 νοικοκυριά στην Ελλάδα, διαθέσιμα από
το \en{University of Strathclyde}.
\bigskip

\textbf{Περιεχόμενο:}
\begin{itemize}
    \item Ηλεκτρολογικές μετρήσεις: ολική ισχύς, τάση, ρεύμα,
    κατανάλωση ανά συσκευή (15 μήνες, ανάλυση 1 λεπτού).
    \item Περιβαλλοντικές μετρήσεις: θερμοκρασία, υγρασία
    (εσωτερικά και εξωτερικά).
    \item Μεταδεδομένα κτιρίου και ανωνυμοποιημένα δημογραφικά
    στοιχεία.
\end{itemize}
\bigskip

\textbf{Κλίμακα:} $\sim$8,1 εκατομμύρια χρονοσφραγίδες ανά τύπο
μέτρησης συνολικά για όλα τα νοικοκυριά. Στην παρούσα υλοποίηση
φορτώθηκε δείγμα 10 εγγραφών από το νοικοκυριό \en{House01}, ώστε
να επιδειχθεί η διαδικασία.
\end{frame}

% ─── Διαφάνεια 17: Σχεδίαση μοντέλου ──────────────────────────────
\begin{frame}[fragile]{Σχεδίαση Μοντέλου \en{NGSI-LD}}
Τα δεδομένα μοντελοποιήθηκαν ως \textbf{πέντε τύποι οντοτήτων} που
αντανακλούν τη φυσική δομή του \en{dataset}:

\selectlanguage{english}
\begin{lstlisting}
Building --hasOccupant--> Person
   |
   +--hasAppliance--> Device (x5)

HouseholdElectricMeasurement --refBuilding--> Building
EnvironmentalMeasurement     --refBuilding--> Building
\end{lstlisting}
\selectlanguage{greek}

Οι στατικές οντότητες (\en{Building}, \en{Person}, \en{Device})
δημιουργούνται μία φορά ανά νοικοκυριό· οι χρονοσειρές
(\en{Measurement}) δημιουργούνται μία φορά ανά χρονοσφραγίδα.
\bigskip

Το σχήμα κληρονομεί τη σημασιολογία των επίσημων
\en{Smart Data Models}, εξασφαλίζοντας διαλειτουργικότητα με άλλα
\en{Data Spaces}.
\end{frame}

%=====================================================================
\section{Υλοποίηση}
%=====================================================================

% ─── Διαφάνεια 18: Βήματα στησίματος ─────────────────────────────
\begin{frame}{Βήματα Υλοποίησης}
Το στήσιμο του \en{Data Space} ακολουθεί επτά διακριτά βήματα,
κάθε ένα από τα οποία αποτελεί προϋπόθεση για το επόμενο:
\bigskip
\begin{enumerate}
    \item Δημιουργία κρυπτογραφικής υποδομής (\en{CA} και
    πιστοποιητικά οντοτήτων).
    \item Παραμετροποίηση \en{iSHARE Satellite}.
    \item Παραμετροποίηση \en{Kong API Gateway}.
    \item Παραμετροποίηση \en{Keyrock}.
    \item Εκκίνηση \en{containers}.
    \item Δημιουργία πολιτικής εξουσιοδότησης.
    \item Φόρτωση δεδομένων.
\end{enumerate}
\bigskip
Όλα τα συστατικά εκτελούνται ως \en{Docker containers}, με
αυτοματοποιημένη ρύθμιση μέσω \en{scripts} \en{Python} και
\en{PowerShell}.
\end{frame}

% ─── Διαφάνεια 19: Πρακτικές προκλήσεις ───────────────────────────
\begin{frame}{Πρακτικές Προκλήσεις}
Η υλοποίηση ανέδειξε αρκετές μη προφανείς προκλήσεις:
\bigskip
\begin{itemize}
    \item \textbf{Μορφή κλειδιού:} Το \en{Kong} απαιτεί
    \en{PKCS\#1} ενώ τα υπόλοιπα συστατικά \en{PKCS\#8} — απαιτείται
    παράλληλη παραγωγή και των δύο μορφών.

    \item \textbf{Επέκταση \texttt{keyUsage}:} Ο \en{Keyrock}
    απορρίπτει πιστοποιητικά χωρίς την επέκταση
    \texttt{digitalSignature}.

    \item \textbf{\en{UUID} έναντι \en{EORI}:} Οι πολιτικές
    εκχώρησης αποθηκεύονται στον \en{Keyrock} με βάση το εσωτερικό
    \en{UUID} του καταναλωτή, όχι το \en{EORI} του.

    \item \textbf{Σειρά αρχικοποίησης:} Ο καταναλωτής πρέπει να
    συνδεθεί στον \en{Keyrock} τουλάχιστον μία φορά πριν δημιουργηθεί
    η πολιτική, ώστε να εκχωρηθεί \en{UUID}.
\end{itemize}
\end{frame}

%=====================================================================
\section{Ροές End-to-End}
%=====================================================================

% ─── Διαφάνεια 20: Σύνοψη ροών ────────────────────────────────────
\begin{frame}{Τρεις Βασικές Ροές Επικοινωνίας}
Το σύστημα υποστηρίζει τρεις διακριτές ροές, οι οποίες αντιστοιχούν
στα τρία βασικά σενάρια χρήσης ενός \en{Data Space}:
\bigskip
\begin{enumerate}
    \item \textbf{Φόρτωση Δεδομένων} (\en{Data Ingestion}): Ο πάροχος
    δημοσιεύει δεδομένα στο \en{Data Space}.
    \item \textbf{Ανάγνωση από Πάροχο}: Ο πάροχος ανακτά τα δικά του
    δεδομένα — πλήρης πρόσβαση χωρίς πρόσθετο έλεγχο εξουσιοδότησης.
    \item \textbf{Ανάγνωση από Καταναλωτή}: Ο καταναλωτής ανακτά
    δεδομένα μέσω του μηχανισμού εκχώρησης εξουσιοδότησης.
\end{enumerate}
\bigskip
Και στις τρεις περιπτώσεις, κάθε αίτημα διέρχεται από το
\en{Kong API Gateway}, εξασφαλίζοντας κεντρικό έλεγχο πρόσβασης.
\end{frame}

% ─── Διαφάνεια 21: Provider flow ──────────────────────────────────
\begin{frame}{Ροή Ανάγνωσης από τον Πάροχο}
Η απλούστερη ροή — ο πάροχος αποδεικνύει την ταυτότητά του και
αποκτά πρόσβαση χωρίς εξωτερική εξουσιοδότηση.
\bigskip
\begin{itemize}
    \item Δημιουργία \en{iSHARE JWT} υπογεγραμμένο με
    \texttt{client.key}.
    \item Αποστολή \en{HTTP GET} στο \en{Kong} με το \en{JWT} ως
    \en{Bearer token}.
    \item Το \en{Kong} επαληθεύει υπογραφή και αλυσίδα πιστοποιητικών.
    \item Επειδή \texttt{iss = EU.EORI.NLPLEGMA} (τοπικός πάροχος),
    η πρόσβαση επιτρέπεται άμεσα.
    \item Προώθηση στον \en{Orion-LD} και επιστροφή αποτελεσμάτων.
\end{itemize}
\bigskip

% TODO: provider sequence diagram
\begin{center}
\textit{[Εικόνα: \en{Provider Sequence Diagram}]}
\end{center}
\end{frame}

% ─── Διαφάνεια 22: Consumer flow ──────────────────────────────────
\begin{frame}{Ροή Ανάγνωσης από τον Καταναλωτή — Η Πιο Σύνθετη Ροή}
Ο καταναλωτής δεν αρκεί να αποδείξει ποιος είναι· πρέπει επίσης να
αποδείξει ότι ο πάροχος του έχει εκχωρήσει άδεια.
\bigskip
\begin{enumerate}
    \item Ο καταναλωτής δημιουργεί \en{iSHARE JWT}.
    \item Στέλνει το \en{JWT} στον \en{Keyrock} (ροή \en{OAuth 2.0
    Client Credentials}).
    \item Ο \en{Keyrock} επαληθεύει τον καταναλωτή μέσω
    \en{Satellite}, αναζητά πολιτική εκχώρησης, και επιστρέφει
    \en{access token} με ενσωματωμένο \en{delegationEvidence}.
    \item Ο καταναλωτής στέλνει το \en{access token} στο \en{Kong}.
    \item Το \en{Kong} αξιολογεί το \en{delegationEvidence} και
    αποφασίζει αν θα επιτραπεί η πρόσβαση.
\end{enumerate}
\bigskip

% TODO: consumer sequence diagram
\begin{center}
\textit{[Εικόνα: \en{Consumer Sequence Diagram}]}
\end{center}
\end{frame}

% ─── Διαφάνεια 23: Πολιτική εκχώρησης ─────────────────────────────
\begin{frame}[fragile]{Πολιτική Εκχώρησης Εξουσιοδότησης}
Η πολιτική που εκδίδει ο πάροχος για τον καταναλωτή ακολουθεί τη
δομή \en{XACML}, μεταφρασμένη σε \en{JSON}:

\selectlanguage{english}
\begin{lstlisting}
{
  "delegationEvidence": {
    "policyIssuer": "EU.EORI.NLPLEGMA",
    "target": { "accessSubject": "<consumer_UUID>" },
    "policySets": [{
      "policies": [{
        "target": {
          "resource": {
            "type": "...HouseholdElectricMeasurement",
            "actions": ["GET"]
          }
        },
        "rules": [{ "effect": "Permit" }]
      }]
    }]
  }
}
\end{lstlisting}
\selectlanguage{greek}

Η πολιτική επιτρέπει στον καταναλωτή να εκτελεί \texttt{GET} στις
ηλεκτρολογικές μετρήσεις, και τίποτε άλλο.
\end{frame}

%=====================================================================
\section{Επίδειξη}
%=====================================================================

% ─── Διαφάνεια 24: Σενάρια επίδειξης ──────────────────────────────
\begin{frame}{Σενάρια Επίδειξης Ελέγχου Πρόσβασης}
Η πειραματική αξιολόγηση εκτελέστηκε σε πραγματικές συνθήκες
λειτουργίας μέσω τριών διακριτών σεναρίων:
\bigskip
\begin{enumerate}
    \item \textbf{Πρόσβαση χωρίς ταυτοποίηση} — επαληθεύει ότι το
    σύστημα απορρίπτει μη εξουσιοδοτημένα αιτήματα.
    \item \textbf{Πρόσβαση με \en{token} παρόχου} — επαληθεύει
    την πλήρη πρόσβαση του παρόχου στα δεδομένα του.
    \item \textbf{Πρόσβαση με \en{token} καταναλωτή} — επαληθεύει
    τον πλήρη κύκλο εκχώρησης εξουσιοδότησης μέσω \en{Keyrock}.
\end{enumerate}
\bigskip
Τα σενάρια εκτελούνται μέσω \en{scripts} \en{PowerShell} που
αυτοματοποιούν τη δημιουργία \en{tokens} και την αποστολή
αιτημάτων μέσω \en{curl}.
\end{frame}

% ─── Διαφάνεια 25: Αποτελέσματα ───────────────────────────────────
\begin{frame}{Αποτελέσματα Επίδειξης}
\begin{table}
\centering
\small
\begin{tabular}{lcc}
\toprule
\textbf{Σενάριο} & \textbf{Απόκριση} & \textbf{Αποτέλεσμα} \\
\midrule
Χωρίς \en{token} & \en{401} & Απόρριψη \\
Πάροχος, ηλεκτρολογικά & \en{200} & Επιτυχία \\
Πάροχος, περιβαλλοντικά & \en{200} & Επιτυχία \\
Πάροχος, κτίριο & \en{200} & Επιτυχία \\
Καταναλωτής, ηλεκτρολογικά & \en{200} & Επιτυχία (μέσω \en{delegation}) \\
\bottomrule
\end{tabular}
\end{table}
\bigskip

Τα αποτελέσματα επιβεβαιώνουν ότι:
\begin{itemize}
    \item Η ψηφιακή κυριαρχία του παρόχου διατηρείται πλήρως.
    \item Ο καταναλωτής έχει πρόσβαση \textbf{αποκλειστικά} στους
    πόρους που του έχουν εκχωρηθεί ρητά.
    \item Καμία αίτηση χωρίς έγκυρο \en{token} δεν φτάνει στα
    δεδομένα.
\end{itemize}
\end{frame}

%=====================================================================
\section{Συμπεράσματα}
%=====================================================================

% ─── Διαφάνεια 26: Συμπεράσματα ───────────────────────────────────
\begin{frame}{Κύρια Συμπεράσματα}
\begin{itemize}
    \item \textbf{Εφικτότητα:} Η δημιουργία ενός λειτουργικού
    \en{Data Space} είναι εφικτή αποκλειστικά με ανοιχτά εργαλεία
    και πρότυπα, προσιτή σε οργανισμούς κάθε μεγέθους.

    \item \textbf{Αποτελεσματικότητα Εκχώρησης:} Ο μηχανισμός
    \en{iSHARE delegation} υλοποιεί αποτελεσματικά την ψηφιακή
    κυριαρχία — ο πάροχος ορίζει με ακρίβεια ποιος έχει πρόσβαση
    σε τι.

    \item \textbf{Αξία της \en{PEP/PDP} Αρχιτεκτονικής:} Ο
    διαχωρισμός μεταξύ επιβολής και απόφασης πολιτικής καθιστά το
    σύστημα ευέλικτο, επεκτάσιμο και διαφανές.

    \item \textbf{Τεκμηρίωση Προκλήσεων:} Η εργασία αναδεικνύει
    και τεκμηριώνει πρακτικές προκλήσεις που δεν είναι εμφανείς
    από τις θεωρητικές προδιαγραφές.
\end{itemize}
\end{frame}

% ─── Διαφάνεια 27: Συνεισφορά ──────────────────────────────────────
\begin{frame}{Συνεισφορά της Εργασίας}
\begin{itemize}
    \item Σχεδίαση και πλήρης υλοποίηση αρχιτεκτονικής \en{Data
    Space} βάσει \en{i4Trust}/\en{iSHARE}, με πέντε ανοιχτά
    συστατικά λογισμικού.

    \item Αυτοματοποιημένα \en{scripts} στησίματος που επιτρέπουν
    την αναπαραγωγή της υλοποίησης σε λίγα λεπτά.

    \item Τεκμηρίωση πρακτικών προκλήσεων (μορφές κλειδιών, σειρά
    αρχικοποίησης, αντιστοίχιση \en{UUID}/\en{EORI}) — γνώση που
    δεν υπάρχει στις προδιαγραφές.

    \item Πλήρης επίδειξη \en{end-to-end} ροών με πραγματικά
    δεδομένα από το \en{Plegma Dataset}.

    \item Σχήμα δεδομένων \en{NGSI-LD} βασισμένο σε
    \en{Smart Data Models}, επεκτάσιμο σε άλλους τομείς ενέργειας.
\end{itemize}
\end{frame}

% ─── Διαφάνεια 28: Περιορισμοί ─────────────────────────────────────
\begin{frame}{Περιορισμοί}
Η παρούσα υλοποίηση έχει συνειδητούς περιορισμούς που πρέπει να
ληφθούν υπόψη:
\bigskip
\begin{itemize}
    \item \textbf{Τοπικό περιβάλλον:} Εκτελείται σε ένα μηχάνημα
    μέσω \en{Docker Compose}. Δεν υλοποιεί \en{HTTPS}, εξισορρόπηση
    φορτίου, ή μηχανισμούς αποκατάστασης βλαβών.

    \item \textbf{Μόνο ροές \en{M2M}}: Η ροή
    \en{Human-to-Machine} (πρόσβαση μέσω \en{browser}) δεν
    υλοποιείται, αν και υποστηρίζεται από τα συστατικά.

    \item \textbf{Ένα \en{dataset}:} Το σύστημα τροφοδοτείται από
    έναν μόνο πάροχο. Σε πραγματικό \en{Data Space} θα υπήρχαν
    πολλαπλοί πάροχοι με ανεξάρτητες υποδομές.
\end{itemize}
\end{frame}

% ─── Διαφάνεια 29: Μελλοντικές επεκτάσεις ─────────────────────────
\begin{frame}{Μελλοντικές Επεκτάσεις}
Βάσει της εμπειρίας από την παρούσα υλοποίηση, προτείνονται:
\bigskip
\begin{itemize}
    \item \textbf{Ροή \en{Human-to-Machine}:} Πρόσβαση χρηστών μέσω
    \en{browser} με ροή \en{OAuth 2.0 Authorization Code} και
    διεπαφή \en{Keyrock}.

    \item \textbf{Ανάπτυξη σε \en{Kubernetes}:} Μεταφορά σε
    περιβάλλον \en{cloud} με χρήση \en{Helm charts} της
    \en{FIWARE Foundation}.

    \item \textbf{Πολλαπλοί Πάροχοι:} Επέκταση σε πραγματικά
    ομοσπονδιακό \en{Data Space} με ανεξάρτητους οργανισμούς.

    \item \textbf{\en{Verifiable Credentials}:} Μετάβαση στη
    νεότερη αρχιτεκτονική του \en{FIWARE Data Space Connector}
    με χρήση \en{OID4VC}.

    \item \textbf{Ανακάλυψη Δεδομένων:} Προσθήκη μηχανισμού
    αυτόματης αναζήτησης διαθέσιμων \en{datasets} από
    καταναλωτές.
\end{itemize}
\end{frame}

%=====================================================================
% ΤΕΛΟΣ
%=====================================================================

% ─── Διαφάνεια 30: Επίδειξη ──────────────────────────────────────
\begin{frame}{Επίδειξη Συστήματος}
\begin{center}
\vspace{1.5cm}
\Huge \textcolor{NTUABlue}{Live Demo}
\bigskip

\normalsize
\textit{Επίδειξη του κύκλου εκχώρησης πρόσβασης σε λειτουργία.}
\end{center}
\end{frame}

% ─── Διαφάνεια 31: Ερωτήσεις ──────────────────────────────────────
\begin{frame}{Ερωτήσεις}
\begin{center}
\vspace{2cm}
\Huge \textcolor{NTUABlue}{Ευχαριστώ!}
\bigskip\bigskip

\normalsize
Ερωτήσεις και σχόλια
\bigskip\bigskip

\small
\textbf{Σωτήριος Κάκος} \\
Εργαστήριο Δικτύων Υπολογιστών, Ε.Μ.Π.
\end{center}
\end{frame}

% ─── Backup διαφάνειες (για ερωτήσεις) ───────────────────────────

\appendix

\begin{frame}{Επιπλέον Πληροφορίες — \en{Backup Slides}}
\begin{center}
\vspace{2cm}
\Huge \textcolor{NTUABlue}{Παράρτημα}
\bigskip

\normalsize
\textit{Επιπλέον διαφάνειες για ερωτήσεις}
\end{center}
\end{frame}

% ─── Backup 1: kong plugin ────────────────────────────────────────
\begin{frame}{\en{Backup}: Το \en{Plugin} \texttt{ngsi-ishare-policies}}
Η λογική του \en{plugin} για κάθε αίτημα:
\bigskip
\begin{enumerate}
    \item Εξάγει το \en{Bearer token} από την κεφαλίδα
    \texttt{Authorization}.
    \item Αποκωδικοποιεί το \en{JWT} και εξάγει την αλυσίδα
    πιστοποιητικών.
    \item Επαληθεύει την υπογραφή με τη βιβλιοθήκη
    \texttt{lua-resty-jwt}.
    \item Αυθεντικοποιείται στον \en{Satellite} και λαμβάνει τη
    λίστα εμπιστευμένων \en{CA}.
    \item Ελέγχει αν ο εκδότης (\texttt{iss}) είναι ο τοπικός πάροχος.
    \item Αν όχι, ελέγχει το ενσωματωμένο \en{delegationEvidence}.
    \item Αν όλα είναι έγκυρα, προωθεί την αίτηση στον \en{Orion-LD}.
\end{enumerate}
\end{frame}

% ─── Backup 2: Data ingestion ─────────────────────────────────────
\begin{frame}{\en{Backup}: Ροή Φόρτωσης Δεδομένων}
Η φόρτωση δεδομένων από τον πάροχο διέρχεται από τους ίδιους
μηχανισμούς ασφαλείας με την ανάγνωση:
\bigskip
\begin{itemize}
    \item Το \en{script} \texttt{orion\_ld\_loading\_script.py}
    διαβάζει \en{CSV} αρχεία με \en{pandas}.
    \item Για κάθε εγγραφή δημιουργεί νέο \en{iSHARE JWT} (λόγω της
    διάρκειας ισχύος των 30 δευτερολέπτων).
    \item Στέλνει \en{HTTP POST} στο \en{Kong} με το \en{entity} σε
    \en{JSON-LD} μορφή.
    \item Το \en{Kong} επαληθεύει την εξουσιοδότηση και προωθεί στον
    \en{Orion-LD}.
\end{itemize}
\bigskip

% TODO: ingestion sequence diagram
\begin{center}
\textit{[Εικόνα: \en{Data Ingestion Sequence Diagram}]}
\end{center}
\end{frame}

% ─── Backup 3: Εικόνα Plegma ──────────────────────────────────────
\begin{frame}{\en{Backup}: Δομή Οντοτήτων \en{NGSI-LD}}
% TODO: εικόνα με το διάγραμμα των 5 οντοτήτων και των σχέσεων
\begin{center}
\vspace{1cm}
\textit{[Εικόνα: Σχέσεις μεταξύ \en{Building}, \en{Person}, \en{Device}, \en{HouseholdElectricMeasurement}, \en{EnvironmentalMeasurement}]}
\end{center}
\bigskip

Αυτή η οργάνωση επιτρέπει ερωτήματα όπως:
\begin{itemize}
    \item «Όλες οι μετρήσεις από νοικοκυριά με ηλιακούς συλλέκτες»
    \item «Μέση κατανάλωση ψυγείων στο \en{dataset}»
\end{itemize}
χωρίς να χρειάζεται \en{join} ή εξωτερική αναζήτηση.
\end{frame}

\end{document}